\chapter{The interface in detail}

MobileCAD's interface has only four permanent areas: the \textbf{top bar}, the \textbf{feature tree}, the
\textbf{toolstrip and view area}, and the \textbf{element tabs}. Every dialog, panel and hint appears as an
overlay on top of the view area. Understand these four and you understand the whole app.

\screenshot{ui-main-zh.jpg}{The four permanent areas of the main screen: top bar (1), feature tree (2), toolstrip and view area (3), element tabs (4).}

\section{Top bar}

From left to right:

\begin{itemize}
  \item \ui{☰}: the main menu (File, Edit, View, Tools, Macro, Help).
  \item \textbf{Document name}: the name of the current document (\file{.FCStd}); it shows “Unnamed” until
        you save.
  \item \textbf{Subtitle}: \ui{Workbench · based on FreeCAD commit}, for example “Part · based on FreeCAD main
        e95a31c”. It tells you which workbench you are in and which FreeCAD version the engine corresponds to.
  \item \ui{↶ / ↷}: Undo and Redo. A long-press (or reading the hint text) shows the name of the operation
        you are about to undo, for example “Undo Pad”.
\end{itemize}

\section{Element tabs and workbenches}

The row of tabs along the bottom corresponds to FreeCAD's workbenches, but reorganized around the design
workflow:

\begin{center}
\small
\begin{tabularx}{\linewidth}{@{}lX@{}}
\toprule
\textbf{Tab} & \textbf{Contents} \\
\midrule
\ui{Part} & A combined workspace with Part Design, Part, Sketch, Surface and Draft; everyday modeling happens
here \\
\ui{Assembly} & Joints, solving and BOM \\
\ui{Drawing} & TechDraw: produce 2D drawings and export them to PDF/SVG/DXF \\
\ui{Variables} & Spreadsheet: parameter tables, aliases and expressions \\
\ui{More ▾} & Every other workbench: BIM, FEM, CAM, Mesh, Mesh Part, Points, Reverse Engineering, Robot,
OpenSCAD, Inspection, Material, Draft, Part, Part Design, Sketcher, Surface and so on \\
\bottomrule
\end{tabularx}
\end{center}

You can switch workbenches in two ways: tap a tab at the bottom, or choose from the full list under
\menupath{View → Workbench}. After you switch to a workbench from More, the tab shows its name (for example
“FEM ▾”).

\section{Feature tree}

The left side is the \textbf{model tree} (FreeCAD's Tree View). It is more than a list: it is the history and
structure of the model:

\begin{itemize}
  \item \textbf{Hierarchy}: a Body contains sketches and features; an assembly contains components and joints;
        a drawing contains pages and views.
  \item \textbf{Icons}: each object type has its own icon (body, pad, sketch, datum plane, point, hole,
        pattern, …), the same as in desktop FreeCAD.
  \item \textbf{State}: red text means the feature has an error; gray or translucent means hidden; a
        strikethrough means suppressed; a small icon after the name means the object needs recomputing
        (Touched); the \ui{Tip} marker shows how far the body has been computed. A small red exclamation mark
        on the right of a feature row means that feature failed; open its dialog to see the error.
  \item \textbf{Tap}: select the object (it highlights in the view); \textbf{double-tap}: open its dialog (on a
        body, double-tap switches the active body); \textbf{long-press}: context menu.
  \item \ui{◀}: collapses the feature tree (it becomes a 40 dp vertical strip; tap it to expand again). On a
        tablet, a 14 dp drag handle resizes the tree between 160 and 480 dp (240 dp by default). On a phone
        the feature tree is an overlay drawer: open it from the tree button in the top bar and close it by
        tapping the dark backdrop.
  \item The count at the top of the tree (for example “15”) is the number of objects in the document. The tree
        \textbf{does not support drag-and-drop reordering}: the model structure is determined by the feature
        history. To change the order, use “Roll tip to here” and the references inside feature dialogs.
\end{itemize}

\screenshot[0.72]{tablet-partdesign.jpg}{The feature tree and a Part Design model. The tree shows the full modeling history — Pad, Groove, Mirrored, Revolution, DatumPlane, Pocket, PolarPattern and others — and that is what “parametric” means: every step can be changed later.}

The menu that a long-press opens (it varies with the object type):

\begin{center}
\small
\begin{tabularx}{\linewidth}{@{}lX@{}}
\toprule
\textbf{Menu item} & \textbf{What it does} \\
\midrule
\ui{Edit …} / \ui{Edit Attachment} & Open the feature dialog, or edit only its attachment (where the sketch or
datum plane is attached) \\
\ui{Rename} & Change the display name (the internal type stays the same) \\
\ui{Hide / Show} & Toggle visibility (also works on sub-elements) \\
\ui{Suppress / Unsuppress} & Skip the feature in the computation for now, to track down errors or compare
designs \\
\ui{Mark to recompute} & Force the object to update at the next recompute \\
\ui{Recompute object} & Recompute the object now \\
\ui{Roll tip to here} & Move the body's computation endpoint to this feature — design exploration by replaying
history \\
\ui{Move into Body} & Merge an external feature such as a fillet or chamfer into the body's history \\
\ui{Toggle active body} & With several bodies, choose which one new features go into \\
\ui{Delete} & Delete the object (Undo brings it back) \\
\bottomrule
\end{tabularx}
\end{center}

\section{Toolstrip and tool panel}

The toolstrip holds a row of \textbf{toolsets}, for example \ui{Sketch}, \ui{Features}, \ui{Part}, \ui{Surface}
and \ui{Draft}. Each toolset corresponds to one or more FreeCAD toolbars:

\begin{itemize}
  \item Tap a toolset → a \textbf{tool panel} unfolds in the top-left corner of the view area, with the buttons
        grouped by FreeCAD toolbar. A button's name is the FreeCAD command name;
  \item at the top of the panel is \ui{Search all tools}: you can search in English or Chinese (for example
        “pad” or “hole”), and the results are grouped by workbench;
  \item \ui{More tools (N)} at the bottom of the panel folds away or expands the commands of this workbench that
        are not on a toolbar, so every command is reachable;
  \item the right side of the toolstrip shows the \textbf{six most recently used} tools, separated by vertical
        lines, so repeated operations are quick. On a phone the toolstrip scrolls horizontally;
  \item tap the toolset again (or tap empty space in the view) to close the panel.
\end{itemize}

\screenshot{ui-toolpanel-zh.jpg}{The “Features” tool panel: sections follow FreeCAD toolbars (modeling features, …), and both command names and icons come from FreeCAD; the top of the panel searches all tools.}

\section{View area}

\begin{itemize}
  \item \textbf{NaviCube} (top right, 112×108 dp): its six faces (top, bottom, front, back, right, left)
        correspond to the standard views; tap a face center to animate to it. The edge strips and corner blocks
        between the faces give 26 directions; the ▲▼◀▶ around it rotate by 15° each time, and ↶ ↷ by 90°. The
        small ☰ menu at its top right switches \ui{Orthographic/Perspective}, \ui{Isometric} and \ui{Fit all}.
        Drag the cube itself to rotate freely. It always shows the current orientation.
  \item \textbf{Selection}: tap to pick faces, edges and vertices. Selected elements collect in a “selection
        bar” along the bottom of the view area (a label plus ×); tap × to deselect.
  \item \textbf{Floating bars}: the sketch status bar (“Fully constrained ›”), sketch value entry, Draft hints,
        assembly joint hints, the selection filter and so on float at the edges of the view area without
        covering the model.
  \item \textbf{Context menu}: \textbf{two fingers tap} on the view (or right-click with a mouse) opens it, and
        the contents follow the current selection: with a face you get “Sketch / Extrude (Pad) / Plane”, with an
        edge “Fillet / Chamfer”, and with a sketch “Extrude (Pad) / Revolve / Pocket / Edit sketch”.
\end{itemize}

\section{The dialog system}

MobileCAD has five kinds of interactive control. Learn to recognize them and you will never be lost:

\begin{description}
  \item[Feature dialog] The equivalent of a desktop FreeCAD task panel: a 44 dp title bar (▾ collapses, ✓
        confirms, ✗ cancels, and the two buttons sit apart so that you do not hit the wrong one), an error strip
        below it, property rows in the middle, and \ui{More options (N)}/\ui{Fewer options} at the bottom. Rows
        follow the order of the matching FreeCAD task panel (for a pad, for example: profile / type / length /
        length 2 / …), with the main rows first and the rest folded away. Row types include: number (Numpad),
        enumeration (up to 4 entries as segmented buttons, more as a drop-down), boolean (switch), link (shows
        “Pad · Edge3, Edge7”; tap to enter picking mode with the \ui{Tap geometry} prompt), color (swatch +
        11-color palette + custom), text/list/file, edge table (a radius per edge), and vector/placement (one
        row per component). \textbf{Every change is an undoable transaction and the model previews live}; ✓
        commits, ✗ rolls back to the state when the dialog opened.
  \item[Task panel] An interface built by FreeCAD's own Python commands (BIM, CAM, FEM, OpenSCAD, TechDraw and
        others) and rendered as the same touch controls: group boxes, number fields, drop-downs, lists, tables,
        checkboxes and radio buttons, and buttons. At the bottom are \ui{OK/Cancel/Apply/Close}, as FreeCAD
        defines them.
  \item[Chooser] A title plus a number of options (pick one or several). Used for “choose a plane”, “choose a
        joint type”, “choose the views to display” and so on.
  \item[Numpad] Slides in from the edge of the screen when a number is needed (on a tablet a 320 dp panel on
        the right, on a phone full width at the bottom): the keys are 7 8 9 / 4 5 6 / 1 2 3 / ± 0 ., the right
        column holds backspace, +, −, × and ÷, and the bottom row holds brackets, C, ✕ and ✓. Units can be
        switched at the top (length mm/cm/m/µm/in/ft, angle °/rad). The entry is a FreeCAD expression, for
        example \file{25 mm} or \file{2 * (10 mm + 5 mm)}. A property value that starts with \file{=} is treated
        as an \textbf{expression} (shown in blue).
  \item[Text input] A single-line or multi-line text box for file names, labels, expressions, macro code and so
        on.
\end{description}

\screenshot{ui-feature-dialog-zh.jpg}{A feature dialog: the title is “Cone”, the properties are Angle / Height / Radius 1 / Radius 2, and the bottom reads “More options (6)”. The arrow at the top left collapses it; ✓ confirms and ✗ cancels at the top right.}

\section{Preferences}

\menupath{Edit → Preferences} (there is also a \ui{Preferences} button on the toolstrip) offers four groups of
settings: unit system, theme, language and update checks. See Section 2.6 for details.

\section{Tablets, phones and headsets}

\begin{itemize}
  \item \textbf{Tablet (width ≥ 600 dp)}: the feature tree stays on the left, the tool panel floats in the
        top-left corner of the view and dialogs float on the left;
  \item \textbf{Phone (width < 600 dp)}: the feature tree becomes a drawer (use the tree button in the top bar,
        or swipe in from the left), panels slide up from the bottom and the toolstrip scrolls horizontally;
  \item \textbf{MobileCAD XR} (Quest and similar headsets): the whole interface is rendered on an 880×620 dp
        panel, the model sits in the room, and the Numpad and on-screen keyboard appear as a second floating
        window; dialogs are drawn inline instead (LocalInlineDialogs in the code).
\end{itemize}

\section{The Back button, Undo and getting stuck}

\begin{itemize}
  \item \textbf{The Back button is Esc}: it closes things from the inside out — Numpad → text dialog → chooser
        → menu → tool panel → Draft tool → assembly joint → task panel → feature dialog → sketch. When
        everything is closed it sends the app to the background \textbf{instead of quitting}, so that a stray
        tap cannot lose your document.
  \item \textbf{Undo/Redo} sits on the right of the top bar. Every feature change and every property edit is a
        separate transaction, and the Undo button's hint names the exact operation.
  \item Unavailable commands (grayed out in menus, 38\% transparent in tool panels) show “<command> is not
        available on mobile yet” when tapped at runtime. They usually belong to desktop window management and
        have no mobile counterpart.
\end{itemize}

\section{Keyboards and input methods}

\begin{itemize}
  \item When text input is needed (expressions, text properties, macros), the on-screen keyboard appears
        automatically and your Android keyboard's word suggestions work.
  \item With a \textbf{hardware keyboard} connected you can type straight into the fields, and physical keys
        are not swallowed by the input method (the app handles Android's IME focus specially for this).
  \item In a single-line field, Enter confirms and hides the keyboard; in a multi-line field, Enter inserts a
        line break.
\end{itemize}
